\clearpage\fancyhead[L]{\small PALLAS.}\SourceLeaf{354}
\begin{center}\textbf{PALLAS.}\par\vspace{4pt}\hrule\vspace{4pt}\fontsize{10}{12}\selectfont Monatliche Correspondenz. Band XXV, S. 389, April 1812.\par\vspace{4pt}\hrule\vspace{4pt}\end{center}

\begin{flushright}Göttingen, den 15. April 1812.\end{flushright}
\par\indent Mit großer Ungeduld erwartete ich dieses Jahr die Wiederkehr der Pallas. Ich wollte sehen, ob sich die Störungsrechnungen diesmal ebenso schön bestätigen würden wie im vorigen Jahr. Mondlicht und schlechtes Wetter ließen mich Pallas aber erst am 4. April sehen. An diesem Abend war es nur kurze Zeit klar, und ich bekam Pallas nur zu Gesicht. Erst am 9. April gelang mir eine Beobachtung, bei der die Deklination allerdings nur geschätzt wurde.

\begin{center}\fontsize{10}{12}\selectfont\setlength{\tabcolsep}{0pt}\renewcommand{\arraystretch}{1.15}\begin{tabular}{l}
1812 April 9\quad $11^{\mathrm u}\;32^{\mathrm m}\;27^{\mathrm s}$ Rektasz. $\PallasSymbol\;269^\circ\;6'\;6''7$.\quad Dekl. $16^\circ\;38'$
\end{tabular}\end{center}

\par\indent Diese Beobachtung zeigt wenigstens, dass die von Herrn \textsc{Nicolai} berechnete Ephemeride vollkommen zutrifft. Es ist sehr zu wünschen, dass man die Opposition am 10. Juni nicht versäumt. Wegen der Lichtschwäche des Planeten wird die Beobachtung mit festen Instrumenten jedoch schwierig sein. Jetzt hat er nur die 11. Größe.

\par\indent Die erste Berechnung der periodischen und säkularen Störungen durch Jupiter habe ich jetzt ganz abgeschlossen. Die Zahl aller entwickelten Gleichungen liegt bei knapp 400. Die größte, die von $5\,\JupiterSymbol-2\,\PallasSymbol$ abhängt, beträgt ungefähr 56 Minuten. Einen einzigen Ort der Pallas zu berechnen, kostet mehrere Tage Arbeit. Einige Ergebnisse dieser Arbeit sind von großem Interesse.

\par\indent Von der gestrigen Bedeckung von $\alpha$ Tauri war hier wegen des bedeckten Himmels nichts zu sehen.
